\documentclass[11pt]{article}

\usepackage[margin=1in]{geometry}
\usepackage{amsmath,amssymb,amsthm}
\usepackage{graphicx}
\usepackage[hidelinks]{hyperref}

\newtheorem{theorem}{Theorem}
\newtheorem{remark}{Remark}

\newcommand{\R}{\mathbb R}

\title{A frontier obstruction for stratified low-dimensional stability}
\author{Agent lane 04}
\date{2026-06-28}

\begin{document}
\maketitle

\begin{abstract}
Alberti, Arroyo and Santacesaria ask whether their stability theory for
low-dimensional inverse problems can be extended beyond fixed-dimensional
manifolds without boundary, motivated by degenerate polygonal parameter
families where the number of parameters varies.  A previous packet proves a
fixed-dimensional boundary/corner extension.  Here we show that the remaining
variable-dimensional stratified upgrade cannot follow from per-stratum
differential injectivity alone: a compact one-dimensional stratum can approach
a zero-dimensional frontier stratum with measurements separated more slowly
than every power of the parameter distance.
\end{abstract}

\section{Source question}

The source paper proves H\"older and Lipschitz stability estimates for nonlinear
inverse problems \(F(x)=y\) when the unknown \(x\) belongs to a compact subset
of a finite-dimensional differentiable manifold \(M\).  One central hypothesis
is that \(F\) is injective and that \(dF_x\) is injective on every tangent
space.  The conclusion includes the following open direction.

\begin{center}
\includegraphics[width=.92\linewidth]{figures/open_question_crop.png}
\end{center}

The source also observes that its main theorems extend readily to a finite
number of pairwise disjoint compact manifold pieces, possibly of different
dimensions.  That extension does not cover compact stratified parameter sets
where a higher-dimensional piece accumulates on a lower-dimensional frontier
piece, as happens in genuine degenerations.

\section{The naive stratified extension}

The tempting full upgrade would say that the source theorem remains true for a
compact finite stratification \(K=\bigcup S_j\) if:
\begin{itemize}
\item each stratum \(S_j\) is a finite-dimensional \(C^1\) manifold;
\item \(F\) is continuous and globally injective on \(K\);
\item \(F|_{S_j}\) is \(C^1\), and \(d(F|_{S_j})_x\) is injective on \(T_xS_j\)
  for every \(x\in S_j\).
\end{itemize}
The example below shows that these conditions alone do not imply even a local
H\"older inverse estimate.

\section{Idea of the proof}

Derivative injectivity is an infinitesimal, same-stratum condition.  It says
that \(F\) does not collapse tangent directions at a fixed point of a fixed
stratum.  It does not say how quickly the image of one stratum approaches the
image of another stratum.

The map \(t\mapsto e^{-1/t}\) on \(t>0\) is strictly increasing and has
positive derivative at every point, so it passes the one-dimensional
differential test.  But the image distance from \(t\) to the frontier point
\(0\) is \(e^{-1/t}\), which is smaller than every positive power of \(t\).
This destroys all H\"older stability estimates across the frontier.

\section{Counterexample}

\begin{theorem}\label{thm:frontier-obstruction}
There exist a compact stratified set \(K\subset\R\) and a continuous injective
map \(F:K\to\R\) such that:
\begin{enumerate}
\item \(K\) is the union of one open one-dimensional stratum and
  zero-dimensional endpoint strata;
\item \(F\) is \(C^\infty\) on each stratum;
\item the differential of \(F\) is injective at every point of the
  one-dimensional stratum, while the zero-dimensional condition is vacuous;
\item for no \(C>0\) and no \(\alpha>0\) does
  \[
    |x-y|\le C|F(x)-F(y)|^\alpha,\qquad x,y\in K,
  \]
  hold.
\end{enumerate}
\end{theorem}

\begin{proof}
Let
\[
  K=[0,1]\subset\R,\qquad S_0=\{0\},\qquad S_1=(0,1),\qquad S_2=\{1\}.
\]
This is a compact set decomposed into one open one-dimensional stratum and two
zero-dimensional endpoint strata.  Define
\[
  F(0)=0,\qquad F(t)=e^{-1/t}\quad(0<t\le 1).
\]
Then \(F\) is continuous on \(K\), strictly increasing, and hence injective.
It is smooth on \(S_1\), and
\[
  F'(t)=t^{-2}e^{-1/t}>0,\qquad 0<t<1.
\]
Thus \(d(F|_{S_1})_t:T_tS_1\simeq\R\to\R\) is injective for every \(t\in S_1\).
On \(S_0\) and \(S_2\), the tangent space is \(\{0\}\), so injectivity of the
differential is vacuous.

Suppose that a H\"older stability estimate held for some \(C>0\) and
\(\alpha>0\).  Applying it to \(x=t\) and \(y=0\) gives
\[
  t\le C e^{-\alpha/t},\qquad 0<t\le 1.
\]
Equivalently,
\[
  \frac{t}{e^{-\alpha/t}}=t e^{\alpha/t}\le C.
\]
But \(t e^{\alpha/t}\to\infty\) as \(t\downarrow0\), a contradiction.  Therefore
no positive H\"older inverse estimate holds on \(K\).
\end{proof}

\section{Relation with the source theorem}

The example does not contradict any theorem in the source paper.  The source
manifold theorem assumes one fixed-dimensional manifold, and its finite-union
extension assumes pairwise disjoint compact pieces \(K_p\subset M_p\).  In a
metric space, finitely many pairwise disjoint compact sets are separated by a
positive distance.  The obstruction here is exactly the missing frontier
regime: the higher-dimensional stratum accumulates at a lower-dimensional
stratum, so the stratum pieces are not separated.

Nor does the example contradict the fixed-dimensional boundary/corner partial
result.  Boundary and corner charts keep one ambient finite dimension and use
closed convex coordinate models.  The present example changes the local
dimension at the frontier.

\section{What a positive theorem would need}

The counterexample isolates a necessary extra ingredient for a genuine
stratified extension.  Besides the source hypotheses on each stratum, one needs
some quantitative compatibility between different strata.  One possible form is
a cross-stratum estimate: for every pair of distinct strata \(S_i,S_j\) that can
approach one another in \(K\), there should be constants \(C_{ij}\) and
\(\alpha_{ij}>0\) such that
\[
  \|x-y\|_X\le C_{ij}\|F(x)-F(y)\|_Y^{\alpha_{ij}},
  \qquad x\in S_i,\ y\in S_j
\]
near their common frontier.  Per-stratum differential injectivity supplies no
such estimate by itself.

\section{Novelty check}

The cheap run indexes were searched for arXiv:2009.00574, the existing
boundary/corner packet, ``stratified frontier'', ``variable-dimensional
degeneracies'', and ``degenerate polygons''.  The only direct hit was the prior
fixed-dimensional partial result for the same source paper.

Web searches on 2026-06-28 for exact or close phrases including
``2009.00574 stratified Holder stability'', ``Inverse problems on
low-dimensional manifolds degenerate polygons'', and ``per-stratum injective
differential Holder stability stratified'' located the source paper but no
separate later answer or this obstruction.

\section{References}

\begin{itemize}
\item G. S. Alberti, A. Arroyo, and M. Santacesaria, \emph{Inverse problems on
  low-dimensional manifolds}, arXiv:2009.00574.
\end{itemize}

\end{document}
